\chapter{The Development of ERP Systems: Architecture, Benefits, and Challenges}
\label{ch:erp-development}

\chapterguide
  {You should understand the idea of functional silos from Chapter 1.}
  {explain how MRP evolved into ERP, describe ERP modularity and shared data, evaluate major benefits and implementation challenges, and identify conditions that make an ERP project more likely to succeed.}

\begin{sourcebasis}
This chapter is based on Tutorial 2 and Chapter 2 of \emph{Concepts in
Enterprise Resource Planning}, with implementation ideas reinforced by Chapter 7.
Tutorial 2 asks students to investigate one ERP benefit and one ERP challenge,
support each with evidence, and teach the result to peers.
\end{sourcebasis}

\section{Why ERP systems emerged}

Business computing did not begin with an integrated platform. Early systems were
built to solve one department's immediate problem: a manufacturing system
tracked materials, an accounting system recorded transactions, a sales system
captured orders. Each was useful alone while leaving the company as a whole hard
to coordinate.

The manufacturing lineage matters most. \term{Material Requirements Planning
(MRP)} calculated what materials were needed and when. Later systems added
capacity, then extended past the factory into sales, purchasing, finance, HR,
and logistics.

\begin{erpfig}[Evolution towards ERP]{Each generation kept the previous scope
and added a question it could not previously answer. ERP is the point at which
the question stops being about the factory and starts being about the
company.}{fig:erp-evolution}
\begin{tikzpicture}[x=1cm,y=1cm]
\draw[draw=UTPMuted!45,line width=0.7pt] (-0.2,0) -- (15.2,0);
\foreach \x in {0.9,4.6,8.3,12.9}{\draw[draw=UTPMuted!45] (\x,-0.12)--(\x,0.12);}
\foreach \x/\era in {0.9/1960s, 4.6/1970s, 8.3/1980s, 12.9/1990s onward}{
  \node[erpnote,below=1.5mm] at (\x,0) {\era};}
\node[erpquiet,minimum width=3.1cm,minimum height=1.35cm,align=center] (a) at (0.9,1.35)
  {\bfseries Inventory control\\[-0.15em]\tiny What do we have\\\tiny in the store?};
\node[erpstep,minimum width=3.1cm,minimum height=1.35cm] (b) at (4.6,1.35)
  {MRP\\[-0.15em]\tiny\mdseries Which materials,\\\tiny\mdseries and by when?};
\node[erpstep,minimum width=3.1cm,minimum height=1.35cm,draw=UTPTeal,fill=UTPPaleTeal] (c) at (8.3,1.35)
  {MRP II\\[-0.15em]\tiny\mdseries Materials \emph{and}\\\tiny\mdseries capacity, cost, labour};
\node[erpstep,minimum width=4.1cm,minimum height=1.35cm,draw=UTPGreen,fill=UTPPaleTeal] (d) at (12.9,1.35)
  {ERP\\[-0.15em]\tiny\mdseries Company-wide processes\\\tiny\mdseries on shared data};
\draw[erpflow] (a)--(b); \draw[erpflow] (b)--(c); \draw[erpflow] (c)--(d);
\node[erpcap,anchor=west,text width=7.0cm,align=left] at (-0.2,3.05)
  {\bfseries scope: one store room \quad$\longrightarrow$\quad one factory};
\node[erpcap,anchor=east,text width=5.0cm,align=right] at (15.2,3.05)
  {\bfseries $\longrightarrow$\quad one enterprise};
\draw[draw=UTPBlue!35,line width=2pt,-{Stealth[length=3mm]}] (-0.2,2.6) -- (15.2,2.6);
\end{tikzpicture}
\end{erpfig}

\section{The core architectural idea: modules around shared data}

ERP products are modular because businesses perform different kinds of work.
Sales needs quotations; manufacturing needs bills of materials; finance needs
invoices and payments. Modularity lets specialised workflows exist while still
drawing on the same underlying records.

\begin{erpfig}[ERP layers]{The architecture in three layers. Users see
specialised modules, but every module reads and writes the same master and
transaction data underneath. Integration lives in the bottom layer, not in the
screens.}{fig:erp-layers}
\begin{tikzpicture}[x=1cm,y=1cm]
% --- layer 1: users ---
\draw[draw=UTPBlue!35,fill=UTPPaleBlue!40,rounded corners=2.5mm] (0,5.05) rectangle (14.6,6.85);
\node[erpcap,anchor=west,color=UTPBlue] at (0.25,6.52) {\bfseries Users and roles};
\foreach \x/\r in {1.95/Sales Manager, 5.25/Purchasing Officer,
                   8.55/Production Planner, 11.85/Accountant}{
  \node[erpdoc,minimum width=2.9cm,minimum height=0.72cm] at (\x,5.75) {\r};}
% --- layer 2: modules ---
\draw[draw=UTPTeal!40,fill=UTPPaleTeal!45,rounded corners=2.5mm] (0,2.45) rectangle (14.6,4.45);
\node[erpcap,anchor=west,color=UTPTeal] at (0.25,4.12) {\bfseries Modules --- specialised workflows};
\foreach \x/\m in {1.35/CRM, 3.35/Sales, 5.35/Purchase, 7.35/Inventory,
                   11.55/Invoicing, 13.5/Employees}{
  \node[erpstep,minimum width=1.85cm,minimum height=0.9cm,draw=UTPTeal,fill=white] at (\x,3.15) {\m};}
\node[erpstep,minimum width=1.85cm,minimum height=0.9cm,draw=UTPTeal,fill=white]
  at (9.45,3.15) {Manu-\\\scriptsize\bfseries facturing};
% --- layer 3: data ---
\draw[draw=UTPNavy!45,fill=UTPNavy!7,rounded corners=2.5mm] (0,-0.35) rectangle (14.6,1.55);
\node[erpcap,anchor=west,color=UTPNavy] at (0.25,1.26) {\bfseries One database --- where the integration actually is};
\node[erpstore,minimum width=6.5cm,minimum height=1.05cm] at (3.75,0.35)
  {\textsc{master data}\\[-0.15em]\tiny\mdseries customers, vendors, products, employees, prices};
\node[erpstore,draw=UTPGreen,fill=UTPPaleTeal,minimum width=6.5cm,minimum height=1.05cm] at (10.85,0.35)
  {\textsc{transaction data}\\[-0.15em]\tiny\mdseries orders, receipts, MOs, deliveries, invoices};
\foreach \x in {1.35,3.35,5.35,7.35,9.45,11.55,13.5}{
  \draw[erplink] (\x,2.45) -- (\x,1.55);
  \draw[erplink] (\x,4.45) -- (\x,5.05);}
\end{tikzpicture}
\end{erpfig}

Two kinds of data are worth separating, because they behave differently and fail
differently.

\begin{erpfig}[Master versus transaction data]{Master data is written rarely and
read constantly; transaction data is written constantly and refers back to the
master records. One bad master record therefore contaminates every transaction
that points at it.}{fig:master-vs-transaction}
\begin{tikzpicture}[x=1cm,y=1cm]
\node[erphead,anchor=west,color=UTPBlue] at (0,3.35) {Master data --- stable, reused};
\foreach \i/\t in {0/Customer: Seri Iskandar Cycling Club,
                   1/Product: Road Bike (Goods{,} tracked),
                   2/Vendor: Ipoh Metal Supply,
                   3/Price: RM 1{,}000}{
  \pgfmathsetmacro{\y}{2.55-\i*0.78}
  \node[erpdoc,anchor=west,minimum width=5.5cm,text width=5.1cm,align=left,
        minimum height=0.62cm] (m\i) at (0,\y) {\t};}
\node[erphead,anchor=west,color=UTPGreen] at (8.6,3.35) {Transaction data --- one per event};
\foreach \i/\t in {0/{SO 001 --- 5 Road Bikes},
                   1/{Delivery WH/OUT/001 --- Done},
                   2/{Invoice INV/2024/001 --- Posted},
                   3/{Payment --- RM 5{,}000 received}}{
  \pgfmathsetmacro{\y}{2.55-\i*0.78}
  \node[erpdoc,anchor=west,minimum width=6.0cm,text width=5.6cm,align=left,
        minimum height=0.62cm,draw=UTPGreen] (t\i) at (8.6,\y) {\t};}
\foreach \i in {0,...,3}{\draw[erplink,-{Stealth[length=1.8mm]}] (t\i.west) -- (m\i.east);}
\node[erpnote,anchor=center,text width=3.1cm] at (7.3,3.35)
  {each transaction \emph{points at} a master record};
\node[erpchipwarn,anchor=west] at (0,-0.75) {change a price here $\ldots$};
\node[erpchipwarn,anchor=west] at (8.6,-0.75) {$\ldots$ and every later document follows};
\end{tikzpicture}
\end{erpfig}

\begin{keyidea}
ERP integration is not achieved by putting every task on one giant screen. It is
achieved by letting specialised modules share master data and create linked
transaction records from the same business events.
\end{keyidea}

\section{Real time means one current version, not zero delay}

In the textbook, \term{real time} means data is current rather than waiting for a
periodic transfer. A confirmed order reaches the functions that need it without
a clerk exporting a file or re-keying it.

This does not mean every metric updates instantly. The useful operational
meaning is that authorised users work from the same current records instead of
separately maintained departmental copies.

\section{Major benefits of ERP}

The benefits described in the source all run through one of four mechanisms.
Naming the mechanism is what turns ``ERP improves efficiency'' into an argument.

\begin{erpfig}[How benefits arise]{Benefits are consequences, not features. Each
chain runs from a structural change, through a change in what people actually
do, to something that can be measured.}{fig:benefit-chain}
\begin{tikzpicture}[x=1cm,y=1cm]
\node[erpcap] at (2.15,3.05) {\bfseries Mechanism};
\node[erpcap] at (7.3,3.05) {\bfseries What changes in the work};
\node[erpcap] at (12.7,3.05) {\bfseries What you can measure};
\foreach \i/\a/\b/\c in {%
  0/{Data integration}/{one record replaces departmental copies}/{\% invoices needing manual correction},
  1/{Process integration}/{one transaction triggers the next step}/{order-to-delivery cycle time},
  2/{Management visibility}/{reports are read, not assembled}/{days to close the books},
  3/{Standardisation}/{units share processes and structures}/{\% transactions on the standard flow}}{
  \pgfmathsetmacro{\y}{2.35-\i*0.9}
  \node[erpstep,minimum width=4.3cm,minimum height=0.72cm] (a\i) at (2.15,\y) {\a};
  \node[erpdoc,minimum width=5.5cm,text width=5.1cm,minimum height=0.72cm] (b\i) at (7.3,\y) {\b};
  \node[erpgood,minimum width=4.8cm,text width=4.4cm,minimum height=0.72cm] (c\i) at (12.7,\y) {\c};
  \draw[erpflow] (a\i)--(b\i); \draw[erpflow] (b\i)--(c\i);}
\end{tikzpicture}
\end{erpfig}

Better customer service, lower inventory, shorter cycle times, and stronger
control are downstream consequences of these four.

\begin{workedexample}
Suppose West End Bicycles keeps the Road Bike price in a salesperson's
spreadsheet and again in Finance's invoice system. A price change must be made
twice and may be missed once. With one integrated product record, the quotation
and the later invoice read the same price.
\end{workedexample}

\section{Why ERP projects are difficult}

ERP affects many departments at once, so implementation is an organisational
change project as much as a technology project. Licence and hardware costs are
the visible part; the costs that sink projects are mostly below the surface.

\begin{erpfig}[The cost of an implementation]{Software licences are the part
budgeted first and the part that matters least to whether the project succeeds.
Everything below the line is process and people work.}{fig:cost-iceberg}
\begin{tikzpicture}[x=1cm,y=1cm]
\draw[draw=UTPBlue!30,fill=UTPPaleBlue!35] (-0.4,1.5) rectangle (14.8,3.1);
\node[erpcap,anchor=west,color=UTPBlue] at (-0.25,2.82) {\bfseries Budgeted first --- visible};
\node[erpstep,minimum width=3.3cm,minimum height=0.8cm] at (2.0,2.05) {Software licences};
\node[erpstep,minimum width=3.3cm,minimum height=0.8cm] at (5.9,2.05) {Hardware / hosting};
\node[erpstep,minimum width=3.3cm,minimum height=0.8cm] at (9.8,2.05) {Consultants};
\node[erpstep,minimum width=2.9cm,minimum height=0.8cm] at (13.2,2.05) {Implementation};
\draw[draw=UTPMuted,line width=1pt,dashed] (-0.4,1.35) -- (14.8,1.35);
\node[erpnote,anchor=east,fill=white,inner sep=1.5pt] at (14.7,1.35) {waterline};
\draw[draw=UTPOrange!35,fill=UTPPaleOrange!45] (-0.4,-2.2) rectangle (14.8,1.2);
\node[erpcap,anchor=west,color=UTPOrange] at (-0.25,0.92) {\bfseries Discovered later --- where projects fail};
\foreach \x/\y/\t in {2.0/0.2/{Business process\\redesign}, 5.9/0.2/{Master-data cleaning\\and migration},
                      9.8/0.2/{Training and\\change management}, 13.2/0.2/{Testing and\\parallel running},
                      2.0/-1.35/{Lost productivity\\during transition}, 5.9/-1.35/{Process-owner\\time},
                      9.8/-1.35/{Customisation and\\future upgrade debt}, 13.2/-1.35/{Post-go-live\\support}}{
  \node[erpwarn,minimum width=3.3cm,minimum height=0.95cm] at (\x,\y) {\t};}
\end{tikzpicture}
\end{erpfig}

Failed implementations are usually traced to process and people problems:
assuming software will repair a bad process, thin planning before configuration,
inadequate training, treating the work as IT-only, weak executive sponsorship,
poor change management, uncontrolled scope or customisation, and difficult
migration of legacy data.

\begin{pitfall}
``Install ERP'' is not a business objective. A project needs a process objective
--- reduce order cycle time, improve inventory accuracy, shorten the financial
close --- or the organisation can spend heavily and reproduce its old problems
inside newer software.
\end{pitfall}

\section{Standard processes versus customisation}

ERP packages embody a particular way of doing business. Every gap between the
package and the company forces a decision, and the cheapest answer is not always
the right one.

\begin{erpfig}[The customisation decision]{Work down the questions before
writing code. Custom features are not forbidden, but each one is paid for again
at every future upgrade.}{fig:customisation-decision}
\begin{tikzpicture}[x=1cm,y=1cm]
\node[erpterm,minimum width=3.1cm] (start) at (1.6,0) {Process gap found};
\node[erpdec,text width=2.4cm] (d1) at (5.5,0) {Does standard\\configuration cover it?};
\node[erpdec,text width=2.5cm] (d2) at (10.2,0) {Is it a real competitive or\\legal need?};
\node[erpgood,minimum width=3.0cm,minimum height=0.9cm] (o1) at (5.5,-2.35) {Configure\\\tiny\mdseries no code, upgrade-safe};
\node[erpgood,minimum width=3.2cm,minimum height=0.9cm] (o2) at (10.2,2.15) {Adopt the standard\\\tiny\mdseries change the process};
\node[erpwarn,minimum width=3.2cm,minimum height=0.9cm] (o3) at (14.0,0) {Customise\\\tiny\mdseries log the upgrade cost};
\draw[erpflow] (start)--(d1);
\draw[erpflow] (d1)--node[erpnote,above,fill=white,inner sep=1pt]{no}(d2);
\draw[erpflow] (d1)--node[erpnote,right,fill=white,inner sep=1pt]{yes}(o1);
\draw[erpflow] (d2)--node[erpnote,right,fill=white,inner sep=1pt]{no}(o2);
\draw[erpflow] (d2)--node[erpnote,above,fill=white,inner sep=1pt]{yes}(o3);
\node[erpnote,text width=4.4cm,anchor=north] at (10.2,-1.6)
  {The useful test: does this process create advantage, or is it merely familiar?};
\end{tikzpicture}
\end{erpfig}

\section{How to judge ERP success}

A successful implementation is not one that merely goes live. It should move
measurable business outcomes: order-to-delivery cycle time, inventory accuracy
and stockout frequency, invoice and payment processing time, days to close,
forecast accuracy, share of transactions completed without re-entry, user
adoption, service levels, and realised cost or revenue change.

Some benefits resist cash valuation. Better visibility, consistent data, and
the capacity to support growth may still be strategically decisive.

\section{Implementation success factors}

\begin{erptable}
\begin{tabularx}{\linewidth}{@{}>{\bfseries\leavevmode\color{ERPNavy}}p{0.25\linewidth}X@{}}
\toprule
\rowcolor{ERPPaleBlue!92}
Success factor & Why it matters \\
\midrule
Executive sponsorship & Gives the project authority, resources, and cross-functional priority. \\
Clear scope and objectives & Prevents uncontrolled expansion and makes benefits measurable. \\
Process ownership & Puts decisions with people who understand the work, not only the technology. \\
Training and communication & Helps users understand both the system and the new process. \\
Change management & Addresses resistance, role changes, and new responsibilities. \\
Data preparation & Poor master data undermines even correctly configured software. \\
Testing and staged readiness & Finds problems before live transactions depend on the system. \\
\bottomrule
\end{tabularx}
\end{erptable}

\section{How Odoo illustrates ERP modularity}

The West End Bicycles guide uses eight Odoo 19 Community applications. Each has
a distinct job; the records are connected.

\begin{odoobox}
Read the Odoo app grid as a modular ERP architecture:
\begin{itemize}
  \item \odooapp{Contacts}: shared customer and vendor identities.
  \item \odooapp{CRM} and \odooapp{Sales}: demand, opportunities, quotations, orders.
  \item \odooapp{Purchase}, \odooapp{Inventory}, \odooapp{Manufacturing}: sourcing, stock, production.
  \item \odooapp{Invoicing}: customer invoices and payments in the Community workflow.
  \item \odooapp{Employees}: departments and the people behind the processes.
\end{itemize}
The same product, customer, vendor, and order data is reused rather than rebuilt
in each app.
\end{odoobox}

\begin{tryit}
Choose one ERP benefit and one implementation challenge from this chapter. For
each write: (1) the mechanism, (2) one operational example at West End Bicycles,
and (3) one metric that would show whether it improved. This mirrors Tutorial 2
while forcing the argument to be measurable.
\end{tryit}

\begin{quickcheck}
\begin{enumerate}
  \item What is the difference between master data and transaction data? Give two Odoo examples of each.
  \item Why is top-management support important in a cross-functional ERP implementation?
  \item Why can copying an inefficient legacy process exactly into an ERP system reduce the expected benefit?
\end{enumerate}
\end{quickcheck}

\begin{chapterrecap}
\begin{itemize}
  \item ERP grew out of inventory control, MRP, and MRP II by widening scope from the store room to the enterprise (\figref{fig:erp-evolution}).
  \item Modules stay specialised but sit on one database of master and transaction data (\figref{fig:erp-layers}).
  \item Benefits arise through integration, visibility, coordination, and standardisation --- each with a measurable consequence (\figref{fig:benefit-chain}).
  \item The costs that decide success are mostly below the waterline: process redesign, data, training, and change management (\figref{fig:cost-iceberg}).
  \item Customise only when configuration cannot meet a real competitive or regulatory need (\figref{fig:customisation-decision}).
\end{itemize}
\end{chapterrecap}
